\section{Diseño y desarrollo de componentes electrónicos (IoT)}

Se diseñará y desarrollará un sistema de adquisición de datos basado en Internet de las Cosas (IoT), compuesto por sensores NDIR (infrarrojos no dispersivos) para la medición de dióxido de carbono (\ch{CO2}) con un rango de detección de 0 a 5000 ppm y resolución de ± 50 ppm (Winsen Sensor, MH-410D) y sensores de gas metano (\ch{CH4}) con un rango de detección de 0 a 10\% en volumen y una resolución de 0.01\% en volumen (Winsen Sensor, MH-441D). Este sistema permitirá monitorear en tiempo real las concentraciones de estos gases, facilitando su análisis y control en aplicaciones como la calidad del aire, seguridad industrial y gestión ambiental.  
Además, se integrarán sensores de temperatura y humedad como el DHT22, con una precisión de ±0.5 °C para la temperatura y ±2 \% para la humedad relativa, sensores de iluminación como el BH1750 con un rango de detección de 1 a 65535 lux, y sensores de pH y nutrientes (NPK) para el monitoreo de las condiciones del sustrato, CWT soil sensor (NPK Type).

Todos estos sensores serán integrados con microcontroladores ESP32, seleccionados por su bajo costo, flexibilidad y capacidad de procesamiento, así como por su conectividad Wi-Fi integrada para la transmisión de datos en tiempo real. Los circuitos electrónicos serán diseñados utilizando software especializado como Eagle, lo que permitirá desarrollar prototipos funcionales y optimizados. Este sistema garantizará la recopilación continua de datos ambientales y de emisiones de GEI durante el proceso de bioconversión, mejorando la precisión y trazabilidad de las mediciones necesarias para optimizar la gestión de residuos orgánicos.

Para la configuración de la red IoT los sensores estarán integrados en dispositivos diseñados específicamente para conectarse a una red IoT mediante protocolos de comunicación eficientes, como LoRa o Wi-Fi. Estos dispositivos permitirán la transmisión de datos en tiempo real hacia una plataforma central, optimizando la cobertura y el consumo energético según las condiciones del entorno.
